\documentclass[12pt,a4paper]{article}

\usepackage[utf8]{inputenc}
\usepackage[T1]{fontenc}
\usepackage[spanish]{babel}
\usepackage[margin=2cm]{geometry}
\usepackage{titlesec}
\usepackage{enumitem}
\usepackage{hyperref}
\usepackage{xcolor}
\usepackage{fontawesome5}
\usepackage{setspace}
\setstretch{1.05}

% Colors
\definecolor{accent}{HTML}{00FF9F}
\definecolor{linkblue}{HTML}{0645AD}
\definecolor{darkbg}{HTML}{0A0A0F}
\definecolor{lighttext}{HTML}{A1A1AA}

\hypersetup{
  colorlinks=true,
  linkcolor=linkblue,
  urlcolor=linkblue
}

% Formatting
\pagestyle{empty}
\setlength{\parindent}{0pt}

\titleformat{\section}
  {\large\bfseries\uppercase}
  {}
  {0pt}
  {}
  [\titlerule]

\titlespacing*{\section}{0pt}{12pt}{6pt}

\begin{document}

% Header
\begin{center}
  {\LARGE\bfseries José Ángel Carretero Montes}\\[6pt]
  {\normalsize Ingeniería Informática \& ADE $\mid$ Data Science}\\[8pt]
  \footnotesize
  \faGlobe\ \href{https://joseangelcarretero.dev}{joseangelcarretero.dev} \hspace{6pt}
  \faEnvelope\ \href{mailto:joseangelcarretero300@gmail.com}{joseangelcarretero300@gmail.com} \hspace{6pt}
  \faLinkedin\ \href{https://www.linkedin.com/in/jos\%C3\%A9-\%C3\%A1ngel-carretero-montes-476008202/}{LinkedIn} \hspace{6pt}
  \faGithub\ \href{https://github.com/0xjoseangel}{0xjoseangel} \hspace{6pt}
  \faTwitter\ \href{https://x.com/0xjoseangel}{X}
\end{center}

% Profile
\section{Perfil}
Estudiante de quinto año del doble grado en Ingeniería Informática (Computación y Sistemas Inteligentes) y ADE en la Universidad de Granada. Especializado en Ciencia de Datos y algoritmos de inteligencia artificial, con experiencia en desarrollo full-stack, automatización de infraestructura y sistemas distribuidos. Orientado a la aplicación de modelos analíticos y técnicas de machine learning en el ámbito financiero y empresarial.

% Skills
\section{Habilidades Técnicas}
\textbf{Lenguajes:} C++, C, Python, SQL, Java, JavaScript, Ruby, Bash \\
\textbf{Frameworks:} Django, React, Vite \\
\textbf{Infraestructura:} Docker, Ansible, Linux, Git, Oracle \\
\textbf{Data Science:} Jupyter, Pandas, NumPy, Matplotlib

% Education
\section{Educación}
\textbf{Doble grado en Ingeniería Informática + ADE} \hfill 2022 -- Presente \\
\textit{Universidad de Granada} \\
Especialidad en Computación y Sistemas Inteligentes

% Languages
\section{Idiomas}
\textbf{Inglés --- C1} (Aptis ESOL, British Council, 2026; Speaking y Writing B2) \hfill \href{https://credentials.britishcouncil.org/28379ba7-48f4-48de-9bec-4a27901aebc4?key=904660ccf4df4f33d680e802a8130aace31f119fae33bb6c3fc9628da4798365}{\faCheckCircle\ Verificar} \\
\textbf{Español} --- Nativo

% Projects
\section{Proyectos}

\textbf{Habla con tu dinero --- Asistente Financiero con IA} \hfill \href{https://github.com/0xjoseangel/reto-unicaja}{\faGithub\ GitHub} \\
\textit{Python, FastAPI, LangGraph, SQLAlchemy, React, Docker}
\begin{itemize}[leftmargin=1.2em,itemsep=1pt,topsep=2pt,parsep=0pt]
  \item Chat financiero (reto Cátedra Unicaja, UGR): el LLM solo propone operaciones; se validan de forma determinista y no se ejecutan sin aprobación humana.
  \item Sin SQL libre: 18 herramientas con consultas SQLAlchemy parametrizadas (sin inyección SQL por diseño); 56 tests.
\end{itemize}

\vspace{6pt}

\textbf{Selección de Carteras con Metaheurísticas} \hfill \href{https://github.com/0xjoseangel/portfolio}{\faGithub\ GitHub} \\
\textit{C++17, OpenMP}
\begin{itemize}[leftmargin=1.2em,itemsep=1pt,topsep=2pt,parsep=0pt]
  \item 14 metaheurísticas (ILS, enfriamiento simulado, genéticos, meméticos\ldots) sobre IBEX 35 y S\&P 500 (457 activos), validadas fuera de muestra.
  \item El memético AM-Best gana en todos los índices; evaluación paralelizada con OpenMP.
\end{itemize}

\vspace{6pt}

\textbf{Danger Theory --- Detección de Crisis Financieras} \hfill \href{https://github.com/0xjoseangel/DANGER-THEORY-DT-}{\faGithub\ GitHub} \\
\textit{Python, Pandas, Algoritmos genéticos}
\begin{itemize}[leftmargin=1.2em,itemsep=1pt,topsep=2pt,parsep=0pt]
  \item Algoritmo genético inmunológico que predice el «Miércoles Negro» (1992) con un día de antelación; validado también en la crisis mexicana de 1994.
  \item Paper académico y libro interactivo en coautoría, supervisados por el Dr. J. L. Sáez Lozano.
\end{itemize}

\end{document}
